% TREŚĆ ROZDZIAŁU: zastąp wskazówki opisem własnego projektu.
% Zachowaj identyfikatory wymagań, aby można było odwołać się do nich w testach.
% Przykładową tabelę można zmienić albo usunąć w całości.
\chapter{Cel, zakres i wymagania}
\label{chap:cel-wymagania}

\section{Cel projektu}
Sformułuj jeden główny, sprawdzalny rezultat projektu oraz ewentualne cele szczegółowe.

\section{Zakres}
Wskaż funkcje, elementy i etapy objęte projektem. Jawnie zapisz istotne elementy pozostające poza zakresem.

\section{Założenia i ograniczenia}
Opisz ograniczenia techniczne, czasowe, sprzętowe, programowe i organizacyjne.

\section{Wymagania}
Nadaj wymaganiom stabilne identyfikatory i kryteria akceptacji. Przykład przedstawiono w \cref{tab:wymagania}.

\begin{table}[htbp]
  \centering
  \caption{Przykładowy rejestr wymagań}
  \label{tab:wymagania}
  \begin{tabularx}{\textwidth}{@{}p{1.5cm}X X@{}}
    \toprule
    ID & Wymaganie & Kryterium akceptacji \\
    \midrule
    F-01 & System realizuje wskazaną funkcję. & Test T-01 zakończony wynikiem pozytywnym. \\
    NF-01 & Czas odpowiedzi nie przekracza ustalonej wartości. & Pomiar w zdefiniowanych warunkach. \\
    \bottomrule
  \end{tabularx}
\end{table}
